%%
%% This is file `sample-sigconf.tex',
%% generated with the docstrip utility.
%%
%% The original source files were:
%%
%% samples.dtx  (with options: `all,proceedings,bibtex,sigconf')
%% 
%% IMPORTANT NOTICE:
%% 
%% For the copyright see the source file.
%% 
%% Any modified versions of this file must be renamed
%% with new filenames distinct from sample-sigconf.tex.
%% 
%% For distribution of the original source see the terms
%% for copying and modification in the file samples.dtx.
%% 
%% This generated file may be distributed as long as the
%% original source files, as listed above, are part of the
%% same distribution. (The sources need not necessarily be
%% in the same archive or directory.)
%%
%%
%% Commands for TeXCount
%TC:macro \cite [option:text,text]
%TC:macro \citep [option:text,text]
%TC:macro \citet [option:text,text]
%TC:envir table 0 1
%TC:envir table* 0 1
%TC:envir tabular [ignore] word
%TC:envir displaymath 0 word
%TC:envir math 0 word
%TC:envir comment 0 0
%%
%%
%% The first command in your LaTeX source must be the \documentclass
%% command.
%%
%% For submission and review of your manuscript please change the
%% command to \documentclass[manuscript, screen, review]{acmart}.
%%
%% When submitting camera ready or to TAPS, please change the command
%% to \documentclass[sigconf]{acmart} or whichever template is required
%% for your publication.
%%
%%
\documentclass[10pt,sigconf,letterpaper,anonymous,nonacm]{acmart}

%%
%% \BibTeX command to typeset BibTeX logo in the docs
\AtBeginDocument{%
  \providecommand\BibTeX{{%
    Bib\TeX}}}

%% Rights management information.  This information is sent to you
%% when you complete the rights form.  These commands have SAMPLE
%% values in them; it is your responsibility as an author to replace
%% the commands and values with those provided to you when you
%% complete the rights form.
\setcopyright{acmlicensed}
\copyrightyear{2018}
\acmYear{2018}
\acmDOI{XXXXXXX.XXXXXXX}

%% These commands are for a PROCEEDINGS abstract or paper.
\acmConference[Conference acronym 'XX]{Make sure to enter the correct
  conference title from your rights confirmation emai}{June 03--05,
  2018}{Woodstock, NY}
%%
%%  Uncomment \acmBooktitle if the title of the proceedings is different
%%  from ``Proceedings of ...''!
%%
%%\acmBooktitle{Woodstock '18: ACM Symposium on Neural Gaze Detection,
%%  June 03--05, 2018, Woodstock, NY}
\acmISBN{978-1-4503-XXXX-X/18/06}


%%
%% Submission ID.
%% Use this when submitting an article to a sponsored event. You'll
%% receive a unique submission ID from the organizers
%% of the event, and this ID should be used as the parameter to this command.
%%\acmSubmissionID{123-A56-BU3}

%%
%% For managing citations, it is recommended to use bibliography
%% files in BibTeX format.
%%
%% You can then either use BibTeX with the ACM-Reference-Format style,
%% or BibLaTeX with the acmnumeric or acmauthoryear sytles, that include
%% support for advanced citation of software artefact from the
%% biblatex-software package, also separately available on CTAN.
%%
%% Look at the sample-*-biblatex.tex files for templates showcasing
%% the biblatex styles.
%%

%%
%% The majority of ACM publications use numbered citations and
%% references.  The command \citestyle{authoryear} switches to the
%% "author year" style.
%%
%% If you are preparing content for an event
%% sponsored by ACM SIGGRAPH, you must use the "author year" style of
%% citations and references.
%% Uncommenting
%% the next command will enable that style.
%%\citestyle{acmauthoryear}

\usepackage[printonlyused]{acronym}
\usepackage{subcaption}
\usepackage{soul}
\usepackage{amssymb}
\usepackage{ragged2e}

\newcommand{\mf}[1]{\textcolor{red}{[Marco -- #1]}}
\newcommand{\om}[1]{\textcolor{purple}{[Orlando -- #1]}}
\newcommand{\js}[1]{\textcolor{blue}{[Jose -- #1]}}
\newcommand{\ael}[1]{\textcolor{orange}{[Andra -- #1]}}
\newcommand{\jo}[1]{\textcolor{teal}{[Jesus -- #1]}}


\acrodef{3GPP}{3rd Generation Partnership Project}
\acrodef{2G}{Second Generation}
\acrodef{3G}{Third Generation}
\acrodef{4G}{Fourth Generation}
\acrodef{5G}{Fifth Generation}
\acrodef{APN}{Access Point Name}
\acrodef{AR}{Augmented reality}
\acrodef{BBU}{Baseband Unit}
\acrodef{BH}{backhaul}
\acrodef{BW}{bandwidth}
\acrodef{BLER}{block error rate}
\acrodef{CA}{Carrier aggregation}
\acrodef{CN}{Core Network}
\acrodef{CID}{Cell ID}
% \acrodef{ccar}{connected car}
\acrodef{C-RAN}{Cloud-RAN}
\acrodef{CV}{Coefficient of Variation}
\acrodef{C-V2X}{Cellular Vehicle-to-Everything}
\acrodef{CU}{Central Unit}
\acrodef{DD}{Decimal Degrees}
\acrodef{DL}{Downlink}
\acrodef{DM}{Demodulation}
\acrodef{DU}{Distributed Unit}
\acrodef{EARFCN}{E-UTRA Absolute Radio Frequency Channel Number}
\acrodef{eMBB}{enhanced Mobile Broadband}%\ac{MBB}}
\acrodef{eNB}{E-UTRAN Node B}
\acrodef{ETSI}{European Telecommunications Standards Institute}
\acrodef{FH}{fronthaul}
\acrodef{Gb}{Gigabit}
\acrodef{GDPR}{General Data Protection Regulation}
\acrodef{GSM}{Global System for Mobile communication}
\acrodef{GSMA}{Global System for Mobile communications Association}
\acrodef{GPS}{Global Positioning System}
\acrodef{HO}{Handover}
\acrodef{HARQ}{hybrid automatic repeat request}
\acrodef{IMEI}{International Mobile Equipment Identity}
\acrodef{IMSI}{International Mobile Subscriber Identity}
\acrodef{IoT}{Internet of Things}
\acrodef{IS}{Indoor Static}
\acrodef{KDE}{Kernel Density Estimate}
\acrodef{KPI}{Key Performance Indicator}
\acrodef{LPWA}{Low-Power Wide-Area}
\acrodef{M2M}{Machine-to-Machine}
% \acrodef{MBB}{Mobile Broadband}
\acrodef{MCC}{Mobile Country Code}
\acrodef{MCS}{modulation and coding scheme}
\acrodef{MEC}{Multi-Access Edge Computing}
\acrodef{M-IoT}{Multimedia IoT}
\acrodef{MH}{midhaul}
\acrodef{ML}{Machine Learning}
\acrodef{mmWave}{millimeter wave}
\acrodef{MME}{Mobility Management Entity}
\acrodef{mMTC}{massive Machine Type Communication}
\acrodef{MNC}{Mobile Network Code}
\acrodef{MNO}{Mobile Network Operator}
\acrodef{MVNO}{Mobile Virtual Network Operator}
\acrodef{MIMO}{Multiple-Input Multiple-Output}
\acrodef{MSC}{Message Sequence Chart}
\acrodef{MSE}{mean-squared-error}
\acrodef{MT-MEC}{multi-tier MEC}
\acrodef{NB-IoT}{Narrowband Internet of Things}
\acrodef{NDA}{Non-Disclosure Agreement}
\acrodef{NG-RAN}{Next Generation-Radio Access Network}
\acrodef{NR}{New Radio}
\acrodef{NR-ARFCN}{New Radio-Absolute Radio Frequency Channel Number}
\acrodef{NSA}{Non-Standalone}
\acrodef{NSPL}{National Statistics Postcode Lookup}
\acrodef{NSA}{Non-Standalone}
\acrodef{LoS}{Line of Sight}
\acrodef{LTE}{Long-Term Evolution}
\acrodef{LTE-M}{LTE for Machine-Type Communications}
\acrodef{OA}{Output Areas}
\acrodef{OD}{Outdoor Driving}
\acrodef{OFDM}{Orthogonal Frequency Division Multiplexing}
\acrodef{ONS}{Office for National Statistics}
\acrodef{OS}{Operating System}
\acrodef{OSD}{Outdoor Static Direct}
\acrodef{OSI}{Outdoor Static Inverse}
\acrodef{OT}{Outdoor Train}
\acrodef{OW}{Outdoor Walking}
\acrodef{PBCH}{Physical Broadcast Channel}
\acrodef{PCI}{Physical Cell ID}
\acrodef{RF}{Radio Frequency}
\acrodef{PSS}{Primary Synchronization Signal}
\acrodef{QCI}{QoS class Identifiers}
\acrodef{RAT}{Radio Access Technology}
\acrodef{RAN}{Radio Access Network}
\acrodef{ROI}{Return on Investment}
\acrodef{RRC}{Radio Resource Control}
\acrodef{RRM}{radio resource management}
\acrodef{RRU}{Remote Radio Unit}
\acrodef{PDCP}{Packet Data Convergence Protocol}
\acrodef{RSRP}{Reference Signal Received Power}
\acrodef{RSRQ}{Reference Signal Received Quality}
\acrodef{RTT}{Round Trip Time}
\acrodef{RU}{Radio Unit}
\acrodef{QoS}{Quality of Service}
\acrodef{QCI}{QoS Class Identifier}
\acrodef{SA}{Standalone}
\acrodef{SIM}{subscriber identification module }
\acrodef{SCS}{Subcarrier Spacing}
\acrodef{SINR}{Signal to Interference and Noise Ratio}
\acrodef{SSB}{Synchronization Signal Block}
\acrodef{SS-RSRP}{Secondary Synchronization Reference Signal Received Power}
\acrodef{SSS}{Secondary Synchronization Signal}
\acrodef{SA}{Standalone}
\acrodef{SGSN}{Serving GPRS Support Node}
\acrodef{SGW}{Serving Gateway}
\acrodef{SDU}{Service Data Unit}
\acrodef{std}{Standard deviation}
\acrodef{PC}{Primary Component}
\acrodef{PCA}{Primary Component Analysis}
\acrodef{PRB}{physical resource block}
\acrodef{TAC}{Type Allocation Code}
\acrodef{TB}{transport block}
\acrodef{TDP}{thermal design power}
\acrodef{TTI}{Transmission Time Interval}
\acrodef{UE}{User Equipment}
\acrodef{UL}{Uplink}
\acrodef{VR}{Virtual Reality}
%\acrodef{V2X}{Vehicle-to-Everything}
\acrodef{UTC}{Universal Time Coordinated}
\acrodef{URLLC}{ultra-reliable low-latency communication}
\acrodef{xDR}{eXtended Detail Record}

%%
%% end of the preamble, start of the body of the document source.
\begin{document}

%%
%% The "title" command has an optional parameter,
%% allowing the author to define a "short title" to be used in page headers.
%\title{Energy-Saving Policies in Base Stations: Assessing Current Strategies and Their Impact}
\title{Energy Savings in the RAN: A/B Testing Fixed Threshold-Based Strategies in Production}

%%
%% The "author" command and its associated commands are used to define
%% the authors and their affiliations.
%% Of note is the shared affiliation of the first two authors, and the
%% "authornote" and "authornotemark" commands
%% used to denote shared contribution to the research.
\author{Ben Trovato}
\authornote{Both authors contributed equally to this research.}
\email{trovato@corporation.com}
\orcid{1234-5678-9012}
\author{G.K.M. Tobin}
\authornotemark[1]
\email{webmaster@marysville-ohio.com}
\affiliation{%
  \institution{Institute for Clarity in Documentation}
  \city{Dublin}
  \state{Ohio}
  \country{USA}
}

%%
%% By default, the full list of authors will be used in the page
%% headers. Often, this list is too long, and will overlap
%% other information printed in the page headers. This command allows
%% the author to define a more concise list
%% of authors' names for this purpose.
\renewcommand{\shortauthors}{Trovato et al.}

%%
%% The abstract is a short summary of the work to be presented in the
%% article.
\begin{abstract}
Reducing energy consumption is a primary goal for the mobile telecommunication industry, with strong environmental and economic sustainability implications. The main target for savings is the Radio Access Network (RAN), which is responsible for 73\% of the total energy costs incurred by operators.
The go-to strategy here is reducing the number of active RAN carriers at locations and times where the demand is reduced and can be served with a smaller wireless capacity.
While the scientific community has been proposing a plethora of complex solutions to switch-off underutilized carriers, production networks largely rely nowadays on static threshold-based strategies that run at individual RAN equipment and are thus very simple to deploy.
However, there are no real-world evaluations about the effectiveness of such approaches in reducing energy consumption and about their impact on the end users.
In this paper, we present an A/B testing of four different fixed threshold-based carrier switch-off methods deployed in production networks serving large geographical regions.
Our study suggests that the capability of such strategies of reducing the energy costs hits a clear barrier if the end users' throughput must be preserved, and provides a strong empirical basis in support of more flexible approaches to save energy at the RAN.
%
%However, special attention and planning are needed to avoid reducing the user quality of experience, such as a lack of service or a significant drop in metrics like throughput. This work focuses on the impact of four different energy-saving approaches deployed on a live network in Europe and how they impact the network's performance and the users' Quality of Service.
\end{abstract}

%%
%% The code below is generated by the tool at http://dl.acm.org/ccs.cfm.
%% Please copy and paste the code instead of the example below.
%%
\begin{CCSXML}
<ccs2012>
   <concept>
       <concept_id>10003033.10003079.10011704</concept_id>
       <concept_desc>Networks~Network measurement</concept_desc>
       <concept_significance>500</concept_significance>
       </concept>
   <concept>
       <concept_id>10003033.10003079.10003082</concept_id>
       <concept_desc>Networks~Network experimentation</concept_desc>
       <concept_significance>500</concept_significance>
       </concept>
   <concept>
       <concept_id>10003033.10003079.10011672</concept_id>
       <concept_desc>Networks~Network performance analysis</concept_desc>
       <concept_significance>500</concept_significance>
       </concept>
 </ccs2012>
\end{CCSXML}

\ccsdesc[500]{Networks~Network measurement}
\ccsdesc[500]{Networks~Network experimentation}
\ccsdesc[500]{Networks~Network performance analysis}


%%
%% Keywords. The author(s) should pick words that accurately describe
%% the work being presented. Separate the keywords with commas.
\keywords{Energy saving, Radio Access Network, carrier switch-off}
%% A "teaser" image appears between the author and affiliation
%% information and the body of the document, and typically spans the
%% page.

\received{20 February 2007}
\received[revised]{12 March 2009}
\received[accepted]{5 June 2009}

%%
%% This command processes the author and affiliation and title
%% information and builds the first part of the formatted document.
\maketitle

\section{Introduction}


%\om{We can start by commenting on how mobile operators are dedicating time and resources to achieve the goal of reducing carbon emission}

The last decade of research, innovation and standardization in 5G have brought advancements to every aspect of cellular systems. 
At the same time, these also translate into a steep increase in terms of the energy consumption by the telco sector. 
With 5G word population coverage reaching 45\% at the end of 2023,\footnote{https://www.ericsson.com/en/reports-and-papers/mobility-report/dataforecasts/network-coverage} the high costs associated with deploying, configuring, and operating 5G equipment concurrently with legacy deployments remains a major hurdle toward realizing the urgent climate target on sustainability and energy efficiency~\cite{huawei_report}.
In particular, reducing the energy and carbon footprint for the telecommunication industry is of paramount importance:
in 2019, a set of operators set the goal of reducing their carbon emissions by 45\% from 2020 to 2030\cite{gsma20205g}.
\ael{are we doing anything around the carbon emissions? I'd tend to keep here only the main metrics that we end up using later on in the paper.}
%In the following years, the number of companies joining this goal has been increasing, totaling up to 70 operators worldwide, setting a new goal of reaching net zero by 2050\cite{gsm2024mobile}.
%For example, Orange has committed to being Net Zero Carbon by 2040 via energy efficiency, renewable energy, the circular economy, or carbon capture~\cite{OrangeCommitment}.
%Nokia is putting special effort into artificial intelligence solutions to reduce the carbon footprint, with estimates of up to 30\% energy savings and less CO2 emission for telecommunications radio networks~\cite{held2021artificial}.

%The 2021's GSMA report\cite{green2021going} mentions that the majority of energy (73\%) is consumed in the RAN. 
The bulk of a mobile network's energy consumption (73\%) corresponds to operating the RAN, and it accounts for ensuring the coverage across thousands of square kilometers, transforming energy into radio waves, and receiving and processing incoming signals.
Approaches such as carrier shutdown, channel shutdown, symbol shutdown that have been emerged ever since the 4G era are helpful solving the 5G energy consumption, but remain insufficient. 
%Though they help the network adapt to the demand in resources
%The remaining consumption distribution comprises data centres (9\%), core (13\%), and operations (5\%).
%\om{Then we move to this is possible to achieve through different ways, like changing the energy sources, and sunsetting other technologies like 2G and 3G, but a substantial amount of work has also been towards increasing the energy efficiency
As end-user connectivity demand soars, operators are looking for a good trade-off between decreasing their networks' energy footprint through adaptive strategies to dynamically power on/off their radio deployments, and their service performance towards users.  
%In other words, as a solution to reduce their network's energy consumption, \acp{MNO} choose to limit selected antennas' operation time. 
%With this approach, operators require special attention and planning to avoid impacting the end-user quality of experience.
% \om{we comment a bit about the difficulty of this approach}
Existing strategies for energy consumption are usually conservative in terms of the breadth of network resources that they impact~\cite{Tan2022}.
This further ensures reduced controlled impact in terms of network performance. 


Energy-saving involves various trade-offs, as was seen in the section on shutdown capabilities. 
One trade-off is between network performance and energy saving. Based on the impact on performance, many energy-saving are lossy as they scale resources and the number of resources is closely related to performance. Therefore, it is difficult to achieve lossless performance\cite{Tan2022}.

It should be noted that most of these works assume that the traffic of the shutting down cells is offloaded to a predefined and fully overlapping neighboring cell. This is far from reality, where the traffic of a shutdown cell cannot be guaranteed to be transferred to a specific desired cell,

Theoretical tools have also been developed to aid network energy efficiency optimization. 

%Previous work has shown that 
Optimal sleeping policies for a single server have been derived only for non-bursty traffic in prior work.\cite{two_thresholds_traffic}.  The optimal sleeping policy is shown to be a two-threshold policy with a wait-and-see feature, i.e., the server would wait a period of time to see if there are any extra arrivals before switching modes.

Even if there are no theoretical results available for more than one server, the findings of this work are in line with today's carrier shutdown solutions, which implement two different sets of conditions for shutdown and reactivation \cite{data_driven_energy23}

\om{and we can close by saying that this approach also contributes to cost reduction and lower energy bills, and we jump to the contributions}

\begin{table}[!t]
    \caption{Energy saving policies tested by the MNO. 
    %\textit{Baseline} represents the policy deployed before the trials.
    } \label{tab:mno-trials}
    \vspace{-2mm}
    \resizebox{1.0\columnwidth}{!}{%
    \centering
    \begin{tabular}{l l c c c c}
    &  & & \multicolumn{2}{c}{\textbf{Thresholds}} &  \\
    \cline{4-5}
        & \textbf{Trial period} & \textbf{Active Time Interval} & \textbf{On} & \textbf{Off} & \textbf{Time To Trigger}  \\
        \hline
         %Baseline & -  & - & - & Feb 8th - Feb 15th \\
         %PoC & - & - & One week & - \\
         \#1 & Feb 22 - Mar 1 & 23h - 6h (night) & 5\% & 10\% & 5 mins \\ 
         \#2 & Mar 1 - Mar 8 & 24h (all day) & 5\%  & 10\% & 5 mins  \\ 
         \#3 & Mar 8 - Mar 15& 24h (all day) & 7\%  & 12\% & 5 mins   \\ 
         \#4 & Mar 15 - Mar 22 & 24h (all day) & 10\%  & 15\% & 5 mins  \\ 
    \end{tabular}
    }
\end{table}

\om{Some of the contributions of this work:}
\begin{itemize}
    \item Evaluate and assess fixed threshold-based energy-saving approaches in a live network.
    \item our findings show that using a single pair of thresholds for all the antennas without accounting for the intrinsic characteristics of the region may fall short
    \item our findings motivate a deeper study and proposal that may be antenna-tailor and or data-driven to increase the efficiency and reduce energy. also, we distill some recommendations that may help further optimize energy-saving policies.
\end{itemize}


\section{Previous works}
\js{Move this section before the conclusions?}

Numerous studies have delved into enhancing energy efficiency within telecommunication networks. This research can be categorized into two main areas: methods for measuring energy consumption and strategies for developing algorithms and methodologies to optimize it. However, to date, there is a gap in understanding the impact of realistic policies in real-world settings. Realistically assessing how different policies perform requires parallel comparisons, akin to A/B testing in the web domain, to evaluate the energy gains or savings of one policy against another.

\textbf{Measuring energy consumption.} The energy consumption of mobile networks is one of the major concerns of the telecom industry. Piovesan et al. \cite{Piovesan2022} proposes a model for characterization of the power consumption of 5G multi-carrier BSs, build on a large data collection campaign. Additionally, Lopes et al. \cite{survey_5g_energy_efficiency} conducted a survey outlining the primary energy efficiency technologies provided by 3GPP NR, together with their main benefits and challenges. Moreover, the community acknowledges the need for sustainability, with efforts underway to quantify the carbon footprint of communication networks \cite{IEA2024}. In the computer systems community, there are existing approaches to quantify and optimize carbon emissions, \cite{CarbonScaler}, by exploit the temporal flexibility inherent to many cloud workloads. In this study, we present utilization and performance data gathered at various times throughout the day. This information can be instrumental in assessing whether base station utilization can indeed benefit from temporal flexibility, a crucial aspect in achieving greater carbon efficiency.

\textbf{Optimizing energy saving policies.} Energy-saving policies involve the dynamic shutdown of base stations during low-traffic periods and their subsequent reactivation as traffic increases.Accurately predicting traffic loads in the near future is crucial for determining the appropriate times to switch base stations on and off. Donevski et al \cite{Donevski2019} propose using dense Neural Network and Recurrent Neural Network paradigms, although their approach lacks empirical validation through real-world trials. Moreover, Domenico et al. \cite{user_transfer_5g} focus on the carrier shutdown approach, enabling base stations to autonomously power down during low-traffic periods by transferring their load to neighboring active base stations. Lastly, it is crucial to acknowledge that such optimization involves the sudden removal and addition of carriers to the network. Therefore, there is a need to develop automated and precise methods for generating configuration parameters for newly added carriers in cellular networks, as suggested by \cite{Mahimkar2021}, such work is orthogonal to the one we are presenting.

 % Section 3: Policies \& Dataset
\begin{table*}[!t]
\caption{Description of the regions under study}
\label{tab:regions}
\centering
\begin{tabular}{lcccccc}
\toprule
\textbf{Region name} & \textbf{Area type} & \textbf{Area $km^2$} & \textbf{Density (antennas/sq. km.)} & \textbf{LTE antennas / site} & \textbf{Capacity / Coverage (\%)}\\

\midrule
Region A & Rural & 4 986.82 & 0.21 & $8.5 \pm 6.28$ & 44.41\% / 55.59\% \\
Region B & Suburban & 189.28 & 5.33 & $10.8 \pm 5.99$ & 47.84\% / 52.16\% \\
\bottomrule
\end{tabular}
\end{table*}

\section{Dataset \& Methodology}
Our study focuses on the cellular network deployment of a top-tier MNO in the UK, with over 50M devices connected. During the period of February 22 to March 22, in 2023, this MNO tested four different energy saving policies using a commercial solution from a RAN equipment vendor (see Table~\ref{tab:mno-trials}). We collected fine-grained KPIs at the antenna level and leveraged this data to make a posterior evaluation of the different policies.

\subsection{RAN energy saving solution}\label{subsec:background-energy-saving}

Figure~\ref{fig:energy_saving_behaviour} illustrates the behavior of the commercial energy saving solution deployed by the MNO. It relies on continuous monitoring of the Physical Resource Block (PRB) utilization on cells, and applies a dynamic switch off policy at the level of \textit{power groups}. These groups include antennas covering the same area over different frequency bands, and has an associated order in which the antennas should be switched off according to the policy. The main configurable parameters are:

\textbf{Active Time Interval:} It sets the time period in which the energy saving policy is active in the network. In our case, Policy 1 is only active over the nighttime (23:00-6:00), while the other Policies 2-4 operate all day (see Table~\ref{tab:mno-trials}).

\textbf{On/Off thresholds:} These values define the levels of PRB utilization that trigger the switch on/off of antennas within a power group. In the MNO trials, each tested policy sets different On/Off thresholds (see Table\ref{tab:mno-trials}) with the goal of evaluating different levels of intensity for energy saving \js{use always intensity instead of aggressiveness?}. When the PRB utilization goes below the Off threshold, the solution estimates the resulting utilization on the remaining active antennas in the power group. This estimate considers that all users attached to the target antenna will move to other antennas of the same power group. If this estimate is below the On threshold, then it starts the switch off process.

\textbf{Time To Trigger:} This parameter sets the hold time from the moment that the On/Off thresholds are exceeded to the moment that the switch On/Off process are actually triggered. This mechanism is intended to prevent erratic ping-pong behaviors due to small load fluctuations around the thresholds. In our case, it is set to 5 minutes in all policies.

To make the switch off process smoother, once a switch off is triggered the antenna power is ramped down before the complete shutdown. This makes that the coverage of the antenna reduces gracefully, so that UEs at the edge start to experience a drop in signal strength and they progressively start to handover to other cells in the same power group. Likewise, the solution applies a contention mechanism that prevents switching off antennas in case there are ongoing emergency calls or queued messages for Earthquake and Tsunami Warning System (ETWS) or Commercial Mobile Alert System (CMAS), thus keeping the network at full capacity under emergency events.
    
\begin{figure}[!t]
    \centering
    \includegraphics[width=1\columnwidth]{images/policies_evaluation/policy_behaviour.pdf}
    \caption{Example of the behavior of the energy saving solution on real-world data. \js{include the configurable parameters: Active Time Interval, on/off thresholds, Time To Trigger; Y label "PRB utilization (\%)"; X label: don't say "hours". It seems the sample period is too long.}.
    \mf{I think that we need this figure (in a better version, using real KPIs over a longer timescale) \emph{right next to Table~\ref{tab:mno-trials}}, so that it is clear what the values in that table refer to}}
    \label{fig:energy_saving_behaviour}
\end{figure}


\subsection{Energy saving trials}\label{subsec:energy-saving-trials}

According to standard deployment strategies, the MNO differentiates its licensed frequency bands into two classes, respectively devoted to optimize \textit{coverage} (low frequency bands) and \textit{capacity} (mid-high frequency bands). From these, capacity antennas are the only ones that implement energy saving policies, while coverage cells typically remain untouched due to their sensitivity to create potential coverage holes in the deployment. Each power group comprises a set of capacity cells along with --~at least~-- one coverage cell. %, which should %$(i)$ provide service to the UEs of the capacity cell, should this one shut down, and $(ii)$ assist the capacity antennas to start the wake up process. Thus, 
Thus, coverage cells not only guarantee service to all users in the area, but also assist inactive capacity cells within the same power group to start the wake up process.  


The studied MNO performed trials of the four policies for one week each (see Table~\ref{tab:mno-trials}), and over two regions identified by specific Tracking Areas. Table~\ref{tab:regions} depicts some details about the two trial regions, which include heterogeneous geographic properties (urban / suburban) and cell deployment densities (cells per sq. km.). Energy saving policies were only tested on LTE antennas, which account for <XX\% \js{Orlando, check this number} of all the antennas deployed in the two regions. Capacity antennas --~which are the ones affected by energy saving policies~-- account for 44\% and 47\% of the tested antennas under Region A and Region B, respectively. Note that the definition of power groups (i.e., coverage and capacity cells under the same coverage area) is a delicate process that may have significant impact on the energy saving results. Nevertheless, the aforementioned trials consider predefined power groups and focus on evaluating the impact of different energy saving configurations across regions.


\subsection{Data feeds}
We collect antenna-level measurements from the three regions of study (Table~\ref{tab:regions}) over the entire period of the trials (Feb 22-Mar 22, 2023). Also, as a reference for our evaluation we collect measurements on the same regions two weeks before the trials (Feb 8-Feb 15, 2023). For this, we rely on two main data feeds that the MNO injects into its BigData measurement platform: 

\textbf{Radio Access Network deployment}:
We collect an inventory of all the antennas deployed in the three regions of study (Table~\ref{tab:regions}). This includes antenna-level information such as the location (lat/lon coordinates), the eNodeB ID, the orientation (azimuth), tilt, antenna manufacturer and model, and the carrier (i.e, frequency channels). This inventory is upgraded on a daily basis reflecting changes in the network topology (e.g., deployment of new cells and decommissioning of older technologies).

\textbf{Radio Access Network KPIs}:
The MNO's measurement infrastructure collects more than 160 KPIs at the antenna level for all the radio access technologies deployed, including from GSM (2G) to the newest 5G NR cells. Measurements are collected in 15-minute intervals and averaged at the antenna level over hourly intervals -- i.e., every 4 measurements -- that are eventually stored and collected in a BigData cluster. In our study, we consider three main KPIs that allow us to evaluate the different energy-saving policies tested: $(i)$ 
antenna availability -- the number of seconds per hour that the antenna was active, $(ii)$ the antenna load -- measured as the number of Physical Resource Blocks (PRB) used for the Downlink --, $(iii)$ and the average user throughput served by the antenna.


\subsection{Evaluation metrics}
\js{Move this to Sec 4??}

A main aspect behind optimizing energy saving in the RAN is evaluating the tradeoff between the energy saved and the impact on the users' performance\cite{Tan2022}. Thus, we take as a reference the following metrics as a proxy to monitor this tradeoff across the different policies and regions tested in the trials: 

\textbf{Total Downtime:} It measures the portion of PRBs utilized across all LTE cells in the regions. It serves as a proxy of the 

\textbf{Portion of cells affected:} It measures the portion of capacity cells where the energy saving policy applied at least 1\% of the time over the 1-week trial period.

\textbf{Overall Physical Resource Block Utilization:} It measures the portion of PRBs utilized across all LTE cells in the tested region. We use it as a proxy of the network congestion in that area.

\textbf{User average throughput:} \js{or delta throughput?}

\js{Motivate how these metrics help measure the tradeoff energy saving vs performance. Do we have a reference to support that the downtime/uptime is a good proxy for estimating energy consumption?}

% \begin{figure}[h]
%     \centering
%     \includegraphics[width=.95\columnwidth]{images/policies_evaluation/avail_kpi_example.pdf}
%     \caption{Example of the availability KPI}
%     \label{fig:avail_kpi}
% \end{figure}

%One of the KPIs reported is the Antenna availability, defined as the number of seconds that the antennas was active in an hour. An example is provided in Figure~\ref{fig:avail_kpi}.

\section{Energy saving policies in action}

As a first step we measure the percentage of downtime that the capacity booster experience during each week of trials. Fig~\ref{fig:percentage_downtime_capacity} show the results for each region. The policy \textit{24h (10/15)} switch off the carries for around 40\% of the time, that represent 67 carrier/hours. Interestingly the \textit{baseline} perform  relative similar to the \textit{24h (5/10)} policy.
In order to explore the reason behaving this metric, we explore how many capacity boosters meet the thresholds condition, show in Fig.~\ref{fig:fraction_antennas_affected}, and the percentage of time that this downtime of this fraction of antennas. Although the trials target around $44\%-47\%$ of the antennas in the regions, not all the antennas meet the switch-off thresholds.

A different view of the dynamic is showed on Fig.~\ref{fig:avail_policies_week} where it represent the mean average availability of the capacity boosters. As can be see, each policy presents a different behaviour.


In this section, we  take as a reference the week of Feb 8 to Feb 15, where there was a less dynamic energy saving policy (hereafter, \textit{Baseline}) deployed by the MNO at the country level. 

we define some sets to study the impact: I) the full TAC, ii) the set of antennas targeted by the policy that actually was affected, iii) the set of antennas targeted by the policy that was NOT affected, and iv) the set of antennas receiving the demand after the antennas switch off.

\begin{figure}
    \centering
    \subfloat[Region A]{\includegraphics[width=.5\columnwidth]{images/policies_evaluation/percentage_downtime_capacity_booster_southeast.pdf}}
    \subfloat[Region B]{\includegraphics[width=.5\columnwidth]{images/policies_evaluation/percentage_downtime_capacity_booster_north_london.pdf}}
    \caption{Capacity booster \% of downtime across policies and regions}
    \label{fig:percentage_downtime_capacity}
\end{figure}


\begin{figure}
    \centering
    \subfloat{
    \includegraphics[width=.8\columnwidth]{images/policies_evaluation/fraction_affected_antennas_and_up_time_legend.pdf}}
    \setcounter{subfigure}{0}
    \subfloat[Region A]{
    \includegraphics[width=.5\columnwidth]{images/policies_evaluation/fraction_affected_antennas_and_up_time_southeast.pdf}}
    \subfloat[Region B]{
    \includegraphics[width=.5\columnwidth]{images/policies_evaluation/fraction_affected_antennas_and_up_time_north_london.pdf}}
    \caption{Percentage of the capacity booster affected by the trials. Not all the capacity boosters meet the thresholds condition.
    % \mf{\st{Fonts too small.} Also, the two 'shades' of color in the bottom are hard to tell apart, can you use brighter/darker colors so as to differentiate better? \om{I change the colors, are they better now?}}
    }
    \label{fig:fraction_antennas_affected}
\end{figure}


\begin{figure*}
    % \Raggedleft
    \subfloat{\includegraphics[width=.48\textwidth]{images/policies_evaluation/avail_kpi_switch_off_legend.pdf}}
    \setcounter{subfigure}{0}
    \centering
    \subfloat[Region A]{
    \includegraphics[width=.95\textwidth]{images/policies_evaluation/avail_kpi_switch_off_set_southeast_mean_std_new.pdf}
    }
    
    \subfloat[Region B]{
    \includegraphics[width=.95\textwidth]{images/policies_evaluation/avail_kpi_switch_off_set_north_london_mean_std_new.pdf}
    }
    \caption{Capacity booster average availability, over the trials.}
    \label{fig:avail_policies_week}
\end{figure*}


\begin{figure}
    \centering
    \subfloat{
    \includegraphics[width=.85\columnwidth]{images/policies_evaluation/fraction_of_available_antennas_legend.pdf}}
    \setcounter{subfigure}{0}
    \subfloat[Region A]{
    \includegraphics[width=.5\columnwidth]{images/policies_evaluation/fraction_of_available_antennas_southeast.pdf}}
    \subfloat[Region B]{
    \includegraphics[width=.5\columnwidth]{images/policies_evaluation/fraction_of_available_antennas_north_london.pdf}}
    \caption{Fraction of the network available over an average day.}
    \label{fig:fraction_network_available}
\end{figure}



\begin{figure}
    \centering    \subfloat{\includegraphics[width=.85\columnwidth]{images/policies_evaluation/delta_kpi_legend.pdf}}
    % \setcounter{subfigure}{0}
    % \subfloat{
    % \includegraphics[width=.48\columnwidth]{images/policies_evaluation/df_tac_policy_delta_t_4g_tti_utilisation_southeast.pdf}}
    % \subfloat{
    % \includegraphics[width=.48\columnwidth]{images/policies_evaluation/df_tac_policy_delta_t_4g_tti_utilisation_north_london.pdf}}
    
    \subfloat[Region A]{
    \includegraphics[width=.48\columnwidth]{images/policies_evaluation/df_tac_policy_delta_t_4g_user_average_throughput_southeast.pdf}}
    \subfloat[Region B]{
    \includegraphics[width=.48\columnwidth]{images/policies_evaluation/df_tac_policy_delta_t_4g_user_average_throughput_north_london.pdf}}
    \caption{Relative change in antenna utilization and user download throughput. 
    % \mf{Do we need to the top plots? They do not add to the story}\om{we can remove it IMO}
    }
    \label{fig:tti_utilization}
\end{figure}

\begin{figure}
    \centering
    \subfloat[Region A]{
    \includegraphics[width=.5\columnwidth]{images/policies_evaluation/percentage_aggressiveness_vs_throughput_boxplot_daylight_southeast.pdf}}
    \subfloat[Region B]{
    \includegraphics[width=.5\columnwidth]{images/policies_evaluation/percentage_aggressiveness_vs_throughput_boxplot_daylight_north_london.pdf}}
    \caption{Policy \textit{Aggressiveness} with respect to average throughput.}
\end{figure}


\begin{figure}
    \centering
    \subfloat[23-6]{
    \includegraphics[width=.485\columnwidth]{images/policies_evaluation/throughput_nearby_antennas_southeast_11_6_scatter.pdf}}
    \subfloat[24h (5/10)]{\includegraphics[width=.485\columnwidth]{images/policies_evaluation/throughput_nearby_antennas_southeast_24x7_1_scatter.pdf}}

    \subfloat[24h (7/12)]{
    \includegraphics[width=.485\columnwidth]{images/policies_evaluation/throughput_nearby_antennas_southeast_24x7_2_scatter.pdf}}
    \subfloat[24h (10/15)]{\includegraphics[width=.485\columnwidth]{images/policies_evaluation/throughput_nearby_antennas_southeast_24x7_3_scatter.pdf}}
    
    \caption{Average throughput of the antennas in the same sector of the capacity booster. Each dot represent a capacity booster, the value in the x-axis represent the down time, while the value in the y-axis is the average change in throughput of the antennas in the same sector with respect the baseline period.}
\end{figure}

\begin{itemize}
    \item Trials evaluation 
    \begin{itemize}
        \item \checkmark Fig.~\ref{fig:fraction_antennas_affected}.Which fraction of the antennas were affected? The difference between the antenna being part of the trial and the antenna being \textbf{actually} affected by the trial.
        
        \item Mean up-time over the trial with error bars
        
        \item \checkmark Fig.~\ref{fig:fraction_network_available} Fraction of the network available over an average day. \om{\st{Color a shade of grey below the coverage line to highlight that you cannot switch off more than this}.}
        
        \item Jesus plot, increase in utilization vs throughput.

        
        \item \checkmark Fig.~\ref{fig:tti_utilization_in_region} Average relative PRB utilization (PRB utilization/Available PRB) \om{Jose asked for a version of this plot only for the antennas affected by the switch off Fig.~\ref{fig:tti_utilization_in_receiving_antennas}.}

        
        \item Plot based on the tti utilization (t\_4g\_tti\_utilisation)
        
        \item how much the antennas are being switched off/on? (number of transitions)
        
        \item \st{Plot based on the tti utilization} (active\_tti\_dl) \om{it is the same as t\_4g\_tti\_utilisation}.
        
        \item \st{we can also show that it doesn't matter that we are picking some KPIs coz they correlate among them}
        
        \item \checkmark Fig.~\ref{fig:user_throughput_in_region} and Fig.~\ref{fig:user_throughput_in_receiving_antennas} User average throughput
    \end{itemize}

    \item shows an example of areas or antennae with a significant decay in QoS while another area/antenna can handle it.
    \item discusses the tradeoff analysis between increase of utilization vs throughput.
    \item Discuss that there is room for improvement, and the policies can be more aggressive (slide 15)
\end{itemize}


\section{Discussion \& Conclusion}

\begin{itemize}
    \item using a single pair of thresholds for all the regions is not a good approach and can be very conservative for some antennas and very disruptive for others
    \item o	The need for more aggressive/data-driven policies
    \item Not only using information about the site, take into account information from nearby switch off, a pesar de ser adaptativo el 24x7 al final se apagan mas por la noche. y quizas si se apaga durante el dia se reduce la emission del carbono
    \item consider the times when the antennas are being switch off
    \begin{itemize}
        \item Directions: show the policies diagram \ref{fig:direction_to_work} Different dimensions to check
    \end{itemize}
\end{itemize}

\begin{figure}
    \centering
    \includegraphics[width=.75\columnwidth]{images/policies_dimensions.png}
    \caption{Different dimension for policies. \mf{We probably do not need this}}
    \label{fig:direction_to_work}
\end{figure}

% \newpage

%%
%% The next two lines define the bibliography style to be used, and
%% the bibliography file.
\bibliographystyle{ACM-Reference-Format}
\bibliography{bib}


%%
%% If your work has an appendix, this is the place to put it.

\end{document}
\endinput
%%
%% End of file `sample-sigconf.tex'.
